\section{Design Decisions}
\label{sec:decisions}

\subsection{Coordination mechanism}

We chose choreography over a broker for the normal path, MongoDB as the only arbiter of who owns each job, and a
periodic reconciliation as the safety net (Section~\ref{sec:coordination}). Three reasons drove the choice:
\begin{itemize}
\item \textbf{No process in the data path.} A free worker takes a book as soon as it is queued, and replicas scale by
  competing for the same queue. No central process calls every service for every book, so none can become the
  bottleneck or the single point of failure of the pipeline.
\item \textbf{Exactly one winner without transactions.} Messages can be lost or repeated, so they cannot decide
  anything. A conditional single-document write gives exactly one winner per transition, without locks or
  transactions, and it is valid across machines because it compares against MongoDB's clock.
\item \textbf{Recovery from persistent state.} Every repair works from \code{book_states}, so a lost message, an expired
  lease or a lost announcement is found and fixed by reading MongoDB. Hazelcast can restart with empty queues.
\end{itemize}

The design has costs, and we accepted them knowingly. It has more moving parts than an orchestrator: leases, the
heartbeat and a six-step reconciliation. It runs two systems, Hazelcast and MongoDB, where polling MongoDB alone would
work. Batching adds wait, because a book waits for its batch to fill or for 30\,s. And the safety net is slow on
purpose: \code{TARANTINO_REQUEUE_SECONDS} is 120\,s, and it must stay above the batch timeout so that \code{control}
never requeues books an indexer is still holding.

\subsection{The five modes measured, and the winner}
\label{sec:winner}

The assignment allows any coordination mechanism, so we also built the four alternatives of Table~\ref{tab:modes} and
measured all five in one campaign (Section~\ref{sec:modes}). The decision rule was written on 5~November, before the
campaign ran:
\begin{enumerate}
\item Correctness is a gate: every applicable row of the failure matrix must recover with the same answers.
\item A difference only counts when \code{pairs.csv} marks it as significant, that is, when the 95\,\% confidence
  interval of the difference between two modes excludes zero.
\item The criteria are then applied in this order: pipeline performance (throughput and time to publication),
  visibility, cost (the CPU of the coordinator and of MongoDB, the containers and the code of the mode), robustness
  (lost work, a broker restart, \code{control} down), and fit for Stage~3.
\item The default stays \mode{hazelcast} in this delivery, and a mode the data favours is proposed for Stage~3.
\end{enumerate}

Table~\ref{tab:decision} applies the rule. Every mode passes the gate: all 165 recovery runs recovered, and
\code{recovery-compare} found the same books and answers for each row in every mode.

\begin{table}[ht]
\centering
\scriptsize
\setlength{\tabcolsep}{3.5pt}
\caption{Decision matrix of the five coordination modes, in the order of the decision rule. Pipeline figures are
the means over 5 repetitions of \run{20261106-012054-coordination-variants-c1376c7}; recovery figures come from the
165 recovery runs and \run{20261106-073911-recovery-delay-c1376c7}. ``N'' is the number of \code{crawler} and
\code{indexer} replicas together, and ``*'' marks a significant difference against \mode{hazelcast}.}
\label{tab:decision}
\begin{tabularx}{\textwidth}{@{}L{2.45cm}YYYYY@{}}
\toprule
Criterion & \mode{hazelcast} & \mode{activemq} & \mode{activemq-persistent} & \mode{polling} & \mode{orchestrator} \\
\midrule
Gate: rows recovered & 14 of 14 & 14 of 14 & 14 of 14 & 11 of 11 & 13 of 13 \\
\addlinespace
Throughput, books/s, N = 1 / 2 / 3 & 28.5 / 46.7 / 58.6 & 28.4 / 45.5 / 57.5 & 28.1 / 43.3 / 55.0* & 27.5 / 47.9 / 55.5 & 27.7 / 44.7 / 34.6* \\
Time to publication p50, s, N = 1 / 2 / 3 & 50.0 / 30.6 / 24.6 & 50.9 / 32.5 / 25.7 & 49.1 / 31.6 / 23.1* & 53.2* / 30.8 / 26.8* & 54.7* / 35.3* / 30.1* \\
\addlinespace
Visibility p50, ms & 38--45 & 39--48 & 38--48 & 2{,}279--2{,}712* & 168--211* \\
\addlinespace
Coordinator CPU, \% of a core, N = 1 / 3 & 3.6 / 8.3 & 5.4* / 12.8* & 7.8* / 16.6* & 3.6 / 7.0* & 10.8* / 12.7* \\
MongoDB CPU, N = 3 & 224\,\% & 219\,\%* & 219\,\% & 247\,\%* & 140\,\%*, doing less work \\
Coordination containers & 1, Hazelcast & 1, ActiveMQ & 1, ActiveMQ with a disk store & None & None \\
Main code of the mode & 192 lines & 328 lines & 328 lines & 104 lines & 506 lines \\
\addlinespace
Lost work (rows 04, 05) & \code{control} requeues after 11.7--11.8\,s & \code{control} requeues after 11.7--11.8\,s & Broker redelivers; reclaimed in 3.8--5.0\,s & Nothing to lose & \code{control} redispatches after 9.7\,s \\
Broker restart (row 12) & Queues lost; requeued after 12\,s with N = 1, 82--97\,s with N = 2--3 & Queues lost; requeued after 11.7\,s & Messages survive; reclaimed in 6.1\,s & No broker & No broker \\
\code{control} down (row 11) & Pipeline goes on & Pipeline goes on & Pipeline goes on & Pipeline goes on & Pipeline stops \\
\addlinespace
Stage~3 fit (argued, not measured) & Native cluster with partition backups, in memory & Network of brokers or a shared store & A shared or replicated disk store; capped at 1{,}460 messages/s & MongoDB load grows with every worker & One dispatcher: a bottleneck and a single point of failure \\
\midrule
Verdict & \textbf{Winner} & Ties on performance, costs more & \textbf{Runner-up}: best recovery & Out on visibility & Out on performance \\
\bottomrule
\end{tabularx}
\end{table}

\textbf{Performance.} \mode{orchestrator} loses: 41\,\% less throughput with 3 replicas, because one thread dispatches
every book one \code{POST} at a time, and 4.7--5.5\,s more time to publication at every N. \mode{polling} keeps the
throughput but takes 2.2--3.2\,s longer to publish with 1 and 3 replicas. \mode{hazelcast} and \mode{activemq} never
differ significantly. \mode{hazelcast} and \mode{activemq-persistent} only differ with 3 replicas, in opposite
directions: \mode{hazelcast} publishes 6\,\% more books per second, and \mode{activemq-persistent} has a 1.6\,s (6\,\%)
lower median time to publication. Its median is lower because persistent sends feed the crawlers more slowly than they
take books. A book waits 1\,ms in \code{books-to-download} instead of 0.8\,s, and slightly less in
\code{books-to-index}, so part of the gain is a later start of each book's clock, not faster work. The same throttling
costs the 6\,\% of throughput. Performance therefore leaves three modes in a tie.

\textbf{Visibility and cost.} \mode{hazelcast}, \mode{activemq} and \mode{activemq-persistent} make a batch searchable
within 38--48\,ms, so visibility does not separate them. Cost does: \mode{hazelcast}'s coordinator uses significantly
less CPU than both ActiveMQ modes at every N, 35\,\% less than \mode{activemq} and 50\,\% less than
\mode{activemq-persistent} with 3 replicas. It also needs no disk store and has the least code of the three. The rule
stops here: \textbf{\mode{hazelcast} wins}, and it remains the mode of Stage~2.

\textbf{What the rule leaves to Stage~3.} Robustness, the criterion after cost, favours \mode{activemq-persistent}.
It is the only mode whose lost work comes back without \code{control}: the broker redelivers what a dead consumer did
not acknowledge, and the book is claimed again 3.8--6.1\,s after the failure, against 11.7\,s for a requeue by
\code{control}. If robustness came before cost, \mode{activemq-persistent} would win. That is why its property, not its
broker, goes into the Stage~3 plan (Section~\ref{sec:stage3}). \code{control} stays necessary in every mode, because it
also releases expired leases and drives rebuilds, and persistence caps the broker at about 1{,}460 messages/s. At the
500--670 books/s that ingestion reaches, with two messages per book, that cap is close.

\subsection{Single-machine deployment}

Everything runs on one host with Docker Compose, as the assignment requires. Addresses are Compose service names, the
datalake is a local volume, and there is one MongoDB server and one broker. Ports, addresses, data paths, replica
counts, thread counts and every timing are environment variables with working defaults, validated at startup.
Replicas are added with \texttt{-{}-scale} and need no restart, because nginx resolves service names on every request
through Docker's DNS. Only nginx publishes ports, and the front-end and operational ports are separate, so a front-end
can never reach an administrative route.

\subsection{API design}

\begin{itemize}
\item \textbf{One search endpoint} with optional filters and pagination, as the assignment asks, plus book details and
  filter values for the front-end use cases.
\item \textbf{Asynchronous state changes}: \code{202} with a \code{Location} to poll, or \code{200} if already done.
  Handlers stay in the millisecond range, cannot exhaust the HTTP pool, and repeated or concurrent requests duplicate
  nothing.
\item \textbf{Consistent reads}: every search reads one immutable snapshot of published batches, so a request never
  mixes two states of the index, and a rebuild switches the whole snapshot at once.
\item \textbf{Problem Details for every error}, including nginx's, with a stable \code{code} and the failing
  \code{parameter}, so a client handles errors without parsing text.
\end{itemize}

\subsection{Storage}

The datalake keeps the Stage~1 \code{batch} layout, which won the Stage~1 comparison: $O(1)$ store and lookup with one
directory per 1{,}000 ids. The inverted index stays in MongoDB, also chosen in Stage~1, but with one postings document
per term and batch instead of one per term (Section~\ref{sec:indexer}). Metadata moves from SQLite to MongoDB: SQLite
allows one writer at a time, which several indexer replicas cannot share, and MongoDB avoids running a third database
for it. Book states also live in MongoDB, because conditional single-document writes are exactly what claims and leases
need. Each generation of the index has its own collections, so a rebuild never disturbs what \code{query} serves.

\subsection{Concurrency model and configuration}

Replicas coordinate only through conditional writes and the broker, and threads inside a process share only immutable
snapshots and \code{java.util.concurrent} structures, so no lock is ever taken. Table~\ref{tab:config} gives the
configuration we keep and the measurement behind each setting.

\begin{table}[ht]
\centering
\small
\caption{Configuration decisions and the evidence behind them (Section~\ref{sec:benchmarks}).}
\label{tab:config}
\begin{tabularx}{\textwidth}{@{}L{3.0cm}L{3.1cm}Y@{}}
\toprule
Setting & Kept & Evidence \\
\midrule
Books per batch & 100; 250 with a larger heap & 250 gives 17\,\% more books/s than 100, but the indexer container peaks at 745--808\,MiB with \code{-Xmx512m}; 500 ran out of memory \\
\code{indexer} replicas & As many as MongoDB allows & Indexing limits the pipeline; 2 and 3 replicas give $\times$1.70 and $\times$2.10, while MongoDB's CPU grows faster than the indexers' \\
\code{crawler} replicas and workers & 1 replica, 8 to 16 workers & Replicas add nothing (521--548 against 504 books/s); 16 workers add 24\,\% over 8; one crawler already downloads about 18 times faster than one indexer publishes \\
\code{query} replicas & 2 & The second replica adds 29\,\% requests/s and cuts latency by 22\,\%; the third adds nothing \\
HTTP threads & Javalin's default, 8 to 250 & 128 threads add nothing over 32 and lose 4\,\% of the requests/s; 8 threads with 32 clients raise latency by 77\,\% \\
Batch announcements & Topic listener plus a 5\,s refresh & A batch is visible in 44\,ms on average with the listener and 2.6\,s without it, at no measurable CPU cost \\
\bottomrule
\end{tabularx}
\end{table}

\subsection{Smaller decisions}

Table~\ref{tab:small-decisions} lists the smaller decisions that shaped the code, each with the alternative it
rejected.

\begin{table}[!ht]
\centering
\small
\caption{Smaller decisions, with the alternative each one rejected.}
\label{tab:small-decisions}
\begin{tabularx}{\textwidth}{@{}L{3.6cm}YL{3.4cm}@{}}
\toprule
Decision & Why & Rejected \\
\midrule
Java~25 & Allowed by the course; virtual threads carry the lease heartbeats cheaply, and records model messages and snapshots & Java~17 of the example \\
Jackson & Javalin's default mapper; one \code{ObjectMapper} for the API, MongoDB documents and broker messages & Gson of the example \\
Generation pointer for rebuilds & Switching it is one atomic write & Renaming three collections, which is not atomic \\
One conditional write per transition, on \code{$$NOW} & Single-document atomicity needs no transactions, and MongoDB's clock is the only clock & Read, decide in Java, write \\
A lease is lost only by a renewal that does not match & A moment without MongoDB should not cost a batch & Losing the lease on any exception \\
Every JVM exits on \code{OutOfMemoryError} & A dead thread behind a healthy \code{/health} is never recovered & Catching \code{Error} in the worker \\
Ten-year Hazelcast connect timeout & The infinite one overflows when the clock steps back & $-1$ \\
JMH with three forks and the protocol in annotations & The error covers differences between JVMs, and every way of launching keeps the protocol & One fork, protocol in a runner \\
\bottomrule
\end{tabularx}
\end{table}
